\chapter{上极限与下极限：习题解答}

\section*{A组}
\addcontentsline{toc}{section}{A组}

\begin{solution}{12.5}
若 \(n=2m\)，尾部最大项是 \(a_n=1+1/(2m)\)，故
\(\beta_{2m}=1+1/(2m)\)；尾部最小的奇数项为
\(a_{2m+1}=-1+1/(2m+1)\)，故
\(\alpha_{2m}=-1+1/(2m+1)\)。
若 \(n=2m-1\)，有
\[
 \beta_{2m-1}=1+\frac1{2m},\qquad
 \alpha_{2m-1}=-1+\frac1{2m-1}.
\]
令 \(m\to\infty\)，得上极限 \(1\)、下极限 \(-1\)。
\end{solution}

\begin{solution}{12.6}
尾部集合 \(A_{n+1}\subseteq A_n\)。子集的上确界不大于母集上确界，故 \(\beta_{n+1}\le\beta_n\)；子集下确界不小于母集下确界，故 \(\alpha_{n+1}\ge\alpha_n\)。
\end{solution}

\begin{solution}{12.7}
\begin{enumerate}[label=(\arabic*)]
 \item 偶数项趋于 \(+\infty\)，奇数项趋于 \(-\infty\)，故上下极限分别为 \(+\infty,-\infty\)。
 \item 指标 \(4k+1\) 的项趋于 \(+\infty\)，\(4k+3\) 的项趋于 \(-\infty\)，故结论相同；偶数项为零不改变端点。
 \item 整列收敛到 \(1\)，故上下极限都为 \(1\)。
\end{enumerate}
\end{solution}

\begin{solution}{12.8}
若有子列趋于 \(+\infty\)，每个尾部上确界均为 \(+\infty\)，所以上极限为 \(+\infty\)。反之，若上极限为 \(+\infty\)，则每个尾部无上界。递归选
\[
 n_{j+1}>n_j,\qquad a_{n_j}>j,
\]
得到趋于 \(+\infty\) 的子列。
\end{solution}

\section*{B组}
\addcontentsline{toc}{section}{B组}

\begin{solution}{12.9}
令 \(s=\limsup a_n\)。取 \(m_1=1\)，由尾部上确界近似性选 \(n_1\ge m_1\)，使
\(a_{n_1}>\beta_{m_1}-1\)。选好 \(n_j\) 后令 \(m_{j+1}=n_j+1\)，再选
\[
 n_{j+1}\ge m_{j+1},\qquad
 a_{n_{j+1}}>\beta_{m_{j+1}}-\frac1{j+1}.
\]
同时 \(a_{n_j}\le\beta_{m_j}\)，而因 \(m_j\to\infty\)，
\(\beta_{m_j}\to s\)。夹逼得 \(a_{n_j}\to s\)，且指标严格递增。
\end{solution}

\begin{solution}{12.10}
若 \(a_{n_j}\to L\)，固定 \(m\)。充分大时 \(n_j\ge m\)，所以
\(\alpha_m\le a_{n_j}\le\beta_m\)。令 \(j\to\infty\)，得
\(\alpha_m\le L\le\beta_m\)。再令 \(m\to\infty\)，即得
\(\liminf a_n\le L\le\limsup a_n\)。
\end{solution}

\begin{solution}{12.11}
对每个尾部，
\[
 \sup_{k\ge n}(-a_k)=-\inf_{k\ge n}a_k.
\]
对 \(n\) 取下确界，得到
\[
 \limsup(-a_n)
 =\inf_n(-\alpha_n)
 =-\sup_n\alpha_n
 =-\liminf a_n.
\]
\end{solution}

\begin{solution}{12.12}
删去有限首项不改变上下极限。此后对每个尾部都有
\[
 \inf_{k\ge n}a_k\le\inf_{k\ge n}b_k,\qquad
 \sup_{k\ge n}a_k\le\sup_{k\ge n}b_k.
\]
分别取相应极限即得结论。
\end{solution}

\begin{solution}{12.13}
对每个 \(n\)，
\[
 \sup_{k\ge n}(a_k+b_k)
 \le \sup_{k\ge n}a_k+\sup_{k\ge n}b_k.
\]
令 \(n\to\infty\) 得所需不等式。严格例取
\(a_n=(-1)^n,b_n=-(-1)^n\)：两上极限均为 \(1\)，和恒为 \(0\)。
\end{solution}

\begin{solution}{12.14}
给定 \(\varepsilon>0\)，最终
\[
 a-\varepsilon+b_n\le a_n+b_n\le a+\varepsilon+b_n.
\]
由上下极限的单调性与平移公式，
\[
 a-\varepsilon+\limsup b_n
 \le\limsup(a_n+b_n)
 \le a+\varepsilon+\limsup b_n.
\]
令 \(\varepsilon\downarrow0\) 得等式。
\end{solution}

\begin{solution}{12.15}
对每个尾部，非负性给出
\[
 \sup_{k\ge n}(a_kb_k)
 \le\left(\sup_{k\ge n}a_k\right)
   \left(\sup_{k\ge n}b_k\right).
\]
两侧均为有限非负数，令 \(n\to\infty\) 并使用乘积极限定理，即得结论。
\end{solution}

\begin{solution}{12.16}
对每个尾部有精确恒等式
\[
 \sup_{k\ge n}\max\{a_k,b_k\}
 =\max\left\{\sup_{k\ge n}a_k,\sup_{k\ge n}b_k\right\}.
\]
对 \(n\to\infty\) 取极限，并使用有限最大值保持极限，得到所求等式。
\end{solution}

\begin{solution}{12.17}
不总成立；取 \(a_n=(-1)^n\)，左端为 \(1\)，右端也是1，此例不否定；改取部分极限为 \(-2,1\) 的交错列，则左端为 \(2\)，而
\(\abs{\limsup a_n}=1\)。对有界实数列，正确关系为
\[
 \limsup\abs{a_n}
 =\max\{\abs{\liminf a_n},\abs{\limsup a_n}\}.
\]
这是因为绝对值在区间 \([\liminf a_n,\limsup a_n]\) 上的最大值出现在端点，且两个端点都由子列实现。
\end{solution}

\begin{solution}{12.18}
设共同有限值为 \(L\)。尾部下确界 \(\alpha_n\uparrow L\)，上确界
\(\beta_n\downarrow L\)。给定 \(\varepsilon>0\)，取 \(N\) 使
\[
 L-\varepsilon<\alpha_N\le\beta_N<L+\varepsilon.
\]
对所有 \(n\ge N\)，\(\alpha_N\le a_n\le\beta_N\)，故
\(\abs{a_n-L}<\varepsilon\)。
\end{solution}

\section*{C组}
\addcontentsline{toc}{section}{C组}

\begin{solution}{12.19}
若 \(a_n\to L\)，则任给 \(\varepsilon>0\)，充分大的尾部全部位于
\((L-\varepsilon/2,L+\varepsilon/2)\)，故振幅不超过 \(\varepsilon\)。反之若 \(\omega_n\to0\)，则
\[
 \limsup a_n-\liminf a_n
 =\lim_n\omega_n=0.
\]
两者有限且相等，由\exerciseref{12.18}原列收敛。
\end{solution}

\begin{solution}{12.21}
令 \(u_n\) 在偶数时为 \(1\)、奇数时为 \(0\)。三种安排为：
\[
 (a_n,b_n)=(u_n,-u_n+1)
\]
时两者上极限均为1且和恒为1；
\[
 (a_n,b_n)=(2u_n-1,1-2u_n)
\]
时两者上极限均为1且和恒为0；
\[
 (a_n,b_n)=(u_n,u_n)
\]
时两上极限均为1且和的上极限为2。
\end{solution}

\begin{solution}{12.23}
不存在无条件等式。取
\[
 a_{2k}=1,\qquad a_{2k-1}=2^{-2k}.
\]
则 \(\limsup a_n^{1/n}=1\)。但从奇项到偶项的比值趋于 \(+\infty\)，所以
\(\limsup a_{n+1}/a_n=+\infty\)。根式判别与比值判别需要额外的规则性条件，不能直接交换。
\end{solution}

\begin{solution}{12.24}
选 \(r\) 使
\[
 \limsup\abs{a_n}<r<1.
\]
由上极限的尾部刻画，存在 \(N\) 使
\[
 \sup_{k\ge N}\abs{a_k}<r.
\]
于是 \(\abs{a_n}\le r\) 对所有 \(n\ge N\) 成立。
\end{solution}

\begin{solution}{12.26}
尾部闭包
\[
 K_n=\overline{\{a_k:k\ge n\}}
\]
递减。对有界列，交集正是部分极限集合；其最大、最小点分别由尾部确界极限给出。集合表达适合证明闭性与嵌套性，尾部确界表达适合运算不等式和定量界，子列表达适合实现端点、构造反例和把量词否定转化为对象。三种语言描述同一尾部结构，但各自压缩的证明步骤不同。
\end{solution}

